%-------------------------
% Резюме в LaTeX
% Формат основан на шаблоне Jake Gutierrez / sb2nov/resume
%-------------------------

\documentclass[a4paper,10pt]{article}

\usepackage[T2A]{fontenc}
\usepackage[utf8]{inputenc}
\usepackage[russian,english]{babel}
\usepackage{latexsym}
\usepackage[empty]{fullpage}
\usepackage{titlesec}
\usepackage[usenames,dvipsnames]{color}
\usepackage{enumitem}
\usepackage[hidelinks]{hyperref}
\usepackage{fancyhdr}
\usepackage{tabularx}
\input{glyphtounicode}

\pagestyle{fancy}
\fancyhf{}
\fancyfoot{}
\renewcommand{\headrulewidth}{0pt}
\renewcommand{\footrulewidth}{0pt}

\addtolength{\oddsidemargin}{-0.55in}
\addtolength{\evensidemargin}{-0.55in}
\addtolength{\textwidth}{1.1in}
\addtolength{\topmargin}{-0.65in}
\addtolength{\textheight}{1.25in}

\urlstyle{same}
\raggedbottom
\raggedright
\setlength{\tabcolsep}{0in}

\titleformat{\section}{
  \vspace{-5pt}\scshape\raggedright\large
}{}{0em}{}[\color{black}\titlerule \vspace{-6pt}]

\pdfgentounicode=1

%-------------------------
% Пользовательские команды
\newcommand{\resumeItem}[1]{
  \item\small{{#1 \vspace{-2pt}}}
}

\newcommand{\resumeSubheading}[4]{
  \vspace{-2pt}\item
    \begin{tabular*}{0.97\textwidth}[t]{l@{\extracolsep{\fill}}r}
      \textbf{#1} & #2 \\
      \textit{\small #3} & \textit{\small #4} \\
    \end{tabular*}\vspace{-7pt}
}

\newcommand{\resumeProjectHeading}[2]{
  \item
    \begin{tabular*}{0.97\textwidth}{l@{\extracolsep{\fill}}r}
      \small #1 & #2 \\
    \end{tabular*}\vspace{-7pt}
}

\renewcommand\labelitemii{$\vcenter{\hbox{\tiny$\bullet$}}$}
\newcommand{\resumeSubHeadingListStart}{\begin{itemize}[leftmargin=0.15in,label={}]}
\newcommand{\resumeSubHeadingListEnd}{\end{itemize}}
\newcommand{\resumeItemListStart}{\begin{itemize}[leftmargin=0.22in]}
\newcommand{\resumeItemListEnd}{\end{itemize}\vspace{-5pt}}

%-------------------------------------------
\begin{document}

%----------ЗАГОЛОВОК----------
\begin{center}
  \textbf{\Huge \scshape Денис Бортников} \\ \vspace{2pt}
  \textbf{\large Junior Fullstack Developer} \\ \vspace{3pt}
  \small
  \href{tel:+79378136099}{\underline{+7 (937) 813-60-99}} $|$
  \href{mailto:dvbortnikov@yandex.ru}{\underline{dvbortnikov@yandex.ru}} $|$
  \href{https://t.me/denis_vag}{\underline{@denis\_vag}} $|$
  \href{https://github.com/MashumaroDev}{\underline{github.com/MashumaroDev}}
\end{center}

%----------О СЕБЕ----------
\section{О себе}
\small{
Студент первого курса магистратуры МГТУ им. Н. Э. Баумана. Разрабатываю web-приложения
полного цикла: проектирую REST API и модели данных, реализую интерфейсы и связываю
frontend с backend. В pet-проекте создал систему управления автосервисом с ролевой
моделью, онлайн-записью, уведомлениями и административной панелью.
}

%----------ОПЫТ РАБОТЫ----------
\section{Опыт работы}
\resumeSubHeadingListStart
  \resumeSubheading
    {Разработчик ПО}{20.03.2026 -- настоящее время}
    {ООО АПК «Малиновка»}{Москва}
    \resumeItemListStart
      \resumeItem{Сопровождаю программное обеспечение и рабочие станции сотрудников: Windows, драйверы, офисные приложения и периферийное оборудование.}
      \resumeItem{Диагностирую технические неисправности, устраняю проблемы совместимости и участвую в поддержке внутренней IT-инфраструктуры предприятия.}
    \resumeItemListEnd
\resumeSubHeadingListEnd

%----------ПРОЕКТЫ----------
\section{Проекты}
\resumeSubHeadingListStart
  \resumeProjectHeading
    {\textbf{DVB --- система управления автосервисом} $|$ \emph{Fullstack pet-проект}}{2026}
    \resumeItemListStart
      \resumeItem{Разработал информационную систему для сети автосервисов: REST API, пользовательский интерфейс, административную панель, тестирование и контейнеризацию.}
      \resumeItem{Декомпозировал предметную область: выделил сущности пользователей, автомобилей, филиалов, подразделений, ремонтных боксов, услуг, механиков, заказов и уведомлений.}
      \resumeItem{Реализовал ролевую модель и рабочие сценарии для клиента, механика, менеджера, администратора филиала и системного администратора.}
      \resumeItem{Разработал онлайн-запись с выбором филиала, услуги, автомобиля, даты и времени; добавил создание и отслеживание заказа в личном кабинете.}
      \resumeItem{Создал CRUD-интерфейсы для сотрудников, филиалов, подразделений и боксов, обработку заказов и просмотр статистики.}
      \resumeItem{Связал frontend и backend: Axios, JWT-аутентификация, обновление сессии, защищённые маршруты и централизованное состояние в Pinia.}
      \resumeItem{Реализовал уведомления в реальном времени на Django Channels и Redis.}
    \resumeItemListEnd

  \resumeProjectHeading
    {\href{https://github.com/MashumaroDev/DVB_Backend}{\textbf{DVB Backend}} $|$ \emph{Python, Django REST, PostgreSQL, Redis, Docker}}{}
    \resumeItemListStart
      \resumeItem{Реализовал REST API, бизнес-логику, собственную модель пользователя, JWT и permissions для разграничения доступа.}
      \resumeItem{Настроил связи Django ORM между заказами, услугами, сотрудниками, филиалами и автомобилями.}
      \resumeItem{Использовал ASGI, WebSocket и Redis для уведомлений; добавил фильтрацию, журналирование, юнит-тесты и нагрузочные сценарии Locust.}
    \resumeItemListEnd

  \resumeProjectHeading
    {\href{https://github.com/MashumaroDev/DVB_Frontend}{\textbf{DVB Frontend}} $|$ \emph{Vue 3, JavaScript, Pinia, Playwright}}{}
    \resumeItemListStart
      \resumeItem{Создал SPA для клиентов и сотрудников: защищённые маршруты, авторизация, онлайн-запись, профиль, автомобили пользователя и история заказов.}
      \resumeItem{Разработал административную панель с CRUD-разделами и визуализацией статистики через Chart.js.}
      \resumeItem{Добавил E2E-сценарии Playwright и их запуск в GitHub Actions.}
    \resumeItemListEnd
\resumeSubHeadingListEnd

%----------ОБРАЗОВАНИЕ----------
\section{Образование}
\resumeSubHeadingListStart
  \resumeSubheading
    {МГТУ им. Н. Э. Баумана}{Москва}
    {Магистратура: Информатика и вычислительная техника, кафедра ИУ6}{09/2026 -- настоящее время}
  \resumeSubheading
    {МГТУ им. Н. Э. Баумана}{Москва}
    {Бакалавриат: Информатика и вычислительная техника, кафедра ИУ6}{09/2021 -- 08/2026}
    \resumeItemListStart
      \resumeItem{Кафедра «Компьютерные системы и сети». Ключевые дисциплины: ООП, компьютерные сети, операционные системы, базы данных, NASM x86/64 и основы компиляции.}
    \resumeItemListEnd
\resumeSubHeadingListEnd

%----------ДОПОЛНИТЕЛЬНОЕ ОБРАЗОВАНИЕ----------
\section{Дополнительное образование}
\resumeSubHeadingListStart
  \resumeSubheading
    {Компьютерные сети}{27.07.2026 -- 10.08.2026}
    {ООО «Академия Дистанционного Образования», повышение квалификации}{72 часа}
\resumeSubHeadingListEnd

%----------ТЕХНИЧЕСКИЕ НАВЫКИ----------
\section{Технические навыки}
\begin{itemize}[leftmargin=0.15in,label={}]
  \small{\item{
    \textbf{Backend}{: Python, Django, Django REST Framework, REST API, JWT, Django Channels} \\
    \textbf{Frontend}{: JavaScript, Vue 3, Vue Router, Pinia, Axios, HTML, CSS, Chart.js} \\
    \textbf{Базы данных}{: PostgreSQL, Redis, Django ORM, миграции} \\
    \textbf{Тестирование}{: юнит-тесты, Playwright, Locust} \\
    \textbf{Инструменты}{: Git, GitHub Actions, Docker, Docker Compose, Vite, Linux} \\
    \textbf{Языки}{: русский --- родной, английский --- B1}
  }}
\end{itemize}

\end{document}
